\definecolor{c1}{RGB}{84,130,53}
% \newcommand{\supportsCovid}[1]{{\color{c1} \textbf{#1}}}
\newcommand{\supportsCovid}[1]{{\color{c1} {#1}}}

\definecolor{c2}{RGB}{218,0,0}
% \newcommand{\refutesCovid}[1]{{\color{c2} \textbf{#1}}}
\newcommand{\refutesCovid}[1]{{\color{c2} {#1}}}

\definecolor{c3}{RGB}{127,127,127}
\newcommand{\wrongContext}[1]{{\color{c3} \textbf{#1}}}



\begin{table*}[t!]
  \setlength{\tabcolsep}{.25em}
  \footnotesize
  \centering

  \begin{tabularx}{\textwidth}{ X }
    \toprule
    \textbf{Claim 1}: Lopinavir / ritonavir have exhibited favorable clinical responses when used as a treatment for coronavirus. \\
    \midrule
    \textbf{Supports}: \dots \supportsCovid{Interestingly, after lopinavir/ritonavir (Kaletra, AbbVie) was administered, $\beta$-coronavirus viral loads significantly decreased and no or little coronavirus titers were observed.} \\
    \midrule
    \textbf{Refutes}: \refutesCovid{The focused drug repurposing of known approved drugs (such as lopinavir/ritonavir) has been reported failed for curing SARS-CoV-2 infected patients.}. It is urgent to generate new chemical entities against this virus \dots \\
    \midrule
    \textbf{Claim 2}: The coronavirus cannot thrive in warmer climates. \\
    \midrule
    \textbf{Supports}: \supportsCovid{...most outbreaks display a pattern of clustering in relatively cool and dry areas...This is because the environment can mediate human-to-human transmission of SARS-CoV-2, and unsuitable climates can cause the virus to destabilize quickly...} \\
    \midrule
    \textbf{Refutes}: \refutesCovid{...significant cases in the coming months are likely to occur in more humid (warmer) climates, irrespective of the climate-dependence of transmission and that summer temperatures will not substrantially limit pandemic growth}. \\
    \bottomrule
  \end{tabularx}

  \caption{Evidence identified by our system as supporting and refuting two claims concerning COVID-19.}
  \label{tbl:covid_example}


\end{table*}
